
%\needspace{4\baselineskip}
\color{uniblue}{\subsection{Защита от повышения частоты (ЗПЧ)}}
\color{black}

\begin{enumerate}[label=\arabic{section}.\arabic{subsection}.\arabic*, labelsep=4pt, leftmargin=0pt, itemindent=57pt]

\item
Защита от повышения частоты предназначена для контроля повышения уровня частоты в сети. ЗПЧ может быть использована для включения выключателей, ранее отключенных от защит ЗСЧ и ЗИЧ. В устройстве защита от повышения частоты представлена независимыми функциями (ступенями) «ЗПЧ 1» и «ЗПЧ 2».

\item
Основные характеристики функций:
\begin{itemize}
    \item погрешность ИО частоты не более 20 мГц;
    \item погрешность ИО частоты в рабочем диапазоне при напряжении входного сигнала 100 В не более 10 мГц; 
    \item время срабатывания и возврата ИО частоты не более 60 мс.
\end{itemize}

\item
ЗПЧ реагирует на превышение частоты выше уставки. Срабатывание ИО функций формирует выходные сигналы «ЗПЧ 1: ИО F>», «ЗПЧ 2: ИО F>» для ЗПЧ 1, ЗПЧ 2 соответственно.

\item
Если до окончания времени таймера выдержки времени на срабатывание частота становилась ниже уставки возврата органа частоты ЗПЧ, то активируется таймер задержки на возврат ИО Т2. Таймер задержки на возврат дает возможность таймеру выдержки времени на срабатывание продолжить работу, даже в случае, если частота кратковременно становилась ниже уставки.

\item
Формирование выходного сигнала «ЗПЧ 1(2): Блок. по U» информирует о блокировке функции по напряжению в результате снижения одного из линейных напряжений ниже уставки. Также блокировка функции осуществляется при наличии входного дискретного сигнала «Блокировка ЗПЧ».

\item
Ввод функций в работу на этапе параметрирования устройства осуществляется своим программным переключателем SGF1 «Ввод\_функции». Информация о введенной функции отображается выходным сигналом «ЗПЧ 1(2): Ввод».

\item
Функции могут быть выведены из работы оперативно путем активации сигнала «ОВ ЗПЧ 1(2), что характеризуется активным состоянием сигнала «ЗПЧ 1(2): Оперативный вывод».

\item
При активном состоянии входного сигнала «ЗПЧ 1 на сигн», «ЗПЧ 2 на сигн» действие соответствующей функции ЗПЧ переводится на сигнал.

\item
Алгоритм работы ЗПЧ 1 выполнен в соответствии с рисунком \ref{zpch:img1}. Алгоритм работы функции ЗПЧ 2 аналогичный. В таблице \ref{zpch:tbl1} приведены параметры, необходимые для настройки ЗПЧ 1. В таблице \ref{zpch1:tbl1} приведены параметры, необходимые для настройки ЗПЧ 2.

\vspace{1mm}
\begin{figure}[H]
\centering
\includegraphics[width=1\textwidth,height=1\textheight,keepaspectratio]{\fbpath/117. ЗПЧ/_latex/img/zpch.pdf}
\captionof{figure}{Функциональная схема логики ЗПЧ 1}\label{zpch:img1}
\end{figure}

\end{enumerate}
